\documentclass[../paper.tex]{subfiles}
\begin{document}
\section{Multiple Commitment Periods}\label{sec:convergence}
So far, we have assumed a single commitment period before the base game, during which players can simultaneously commit. In this section, we extend the model to allow for multiple commitment periods, denoted by $T$. If players can credibly commit to a strategy for a game they will play in the future, it seems natural that they might be able to do so at several points in time leading up to the game. This is particularly relevant if one thinks of these players as AI agents. Fast computational speed could imply that agents have many, effectively infinite, opportunities to commit.
 
From this point forward, let $T$ represent the number of discrete commitment periods. 
Suppose that in the first commitment period, player $1$ commits and player $2$ waits. Then, player $2$ gets to observe $1$'s commitment in the second period, and can best respond to it. 
But if instead both players did not commit in the first period, they observe each other’s inaction and then in the second period simultaneously decide again whether to commit. The base game is played only if neither player commits in any of the $T$ commitment periods. 

Let $A_{1t}$ denote the strategy whereby a player waits until period $t$ to commit to $A_1$, unless the other player has already committed in an earlier period. In that case, the player best responds. Waiting ($W$) is the strategy of not committing in any commitment period, best-responding to any commitment made by the other player, and if the base game is reached playing the MSNE strategy of the base game. 
Let $p_{t}$ denote the probability of a player playing $A_{1t}$ and let $q_t$ denote the probability of a player playing $A_{2t}$. The probability of waiting is denoted as $\sigma(W)=1-\sum_{\tau=1}^{T}(p_\tau+q_\tau)$.  
 
\begin{proposition}\label{theo:convergence}
    As $T\to\infty$, the expected payoff in the unique mixed subgame perfect equilibrium for both players
    \begin{enumerate}
        \item decreases in the Game of Chicken and converges to $\cparam$
        \item increases in the Battle of Sexes and converges to $\min\{\bparam,\cparam\}$
    \end{enumerate}
\end{proposition}
\noindent 
The full proof of proposition \ref{theo:convergence} can be found in the  appendix in section \ref{sec:prooftheo2}. we prove three lemmas to establish that this is indeed the unique mixed subgame perfect equilibrium for all $T$, and to characterize the equilibrium as $T \to \infty$, including closed-form expressions for all $p_t$, $q_t$, and $\sigma(W)$ in terms of $\bparam$, $\cparam$, and $T$. Table \ref{tab:convergence} describes the probabilities of committing in the first period and waiting as $T\to\infty$, as well as the equilibrium payoffs.
\begin{table}[h!]
\centering
\begin{tabular}{C{3.5cm} C{7cm} C{2cm}}  % adjust widths as needed
\toprule
\textbf{Condition} & \textbf{Asymptotic Behavior} & \textbf{Payoff} \\ \midrule
$0<\cparam \leq 1 < \bparam$ & $p_1\to\dfrac{\bparam-\cparam}{\,\bparam},\quad \sigma(W) \to 0$ & $\to \cparam$ \\[2ex]
$1< \cparam < \bparam$ & 
$p_1 \to \dfrac{\bparam-\cparam}{\,\bparam},\quad q_1 \to 0,\quad \sigma(W) \to 0$ & $\to \cparam$ \\[2ex]
$1< \bparam < \cparam$ &
$p_1 \to 0,\quad q_1 \to \dfrac{\cparam-\bparam}{\,\cparam-1},\quad \sigma(W) \to 0$ & $\to \bparam$ \\ \bottomrule
\end{tabular}
\caption{Asymptotic convergence of probabilities as $T \to \infty$.}
\label{tab:convergence}
\end{table}
\newline
To illustrate Proposition \ref{theo:convergence}, recall the example commitment Game of Chicken from the introduction. With a single commitment period, $T = 1$, the mixed-strategy equilibrium involves both players committing to Straight with probability $\tfrac{8}{11}$, waiting with probability $\tfrac{3}{11}$, and receiving an expected payoff of $-\tfrac{5}{11}$. When the number of commitment periods increases to $T = 2$, the mixed equilibrium changes: each player commits to Straight in the first period with probability $\tfrac{96}{129}$, in the second period with probability $\tfrac{24}{129}$ (conditional on no one having committed in the first period), and waits in both periods with probability $\tfrac{9}{129}$. The resulting expected payoff is $-\tfrac{21}{43}$, which is strictly lower than under $T = 1$. As $T$ increases further, the mixed equilibrium payoff further decreases. The probability of committing Straight in the first period increases and converges to $\tfrac{3}{4}$. The expected payoff converges to $-\tfrac{1}{2}$ for both players, which is the payoff of `losing' the Game of Chicken.

This convergence occurs because each additional commitment period hurts both players. As previously argued, introducing one period of commitment transforms the base Game of Chicken into a game with lower payoffs, as illustrated in Table~\ref{tab:prop1-intuition}. Adding another commitment period to this game transforms it into a game with even lower payoffs. In the first commitment period, players choose whether to commit Straight or wait. If both wait, the game proceeds. Under $T = 1$, this would lead to the MSNE of the base game. Under $T = 2$, however, it leads to the $T = 1$ commitment game, which has a lower expected payoff. As a result, the equilibrium payoff is lower when $T=2$. Each additional commitment period worsens outcomes in the Game of Chicken. As $T\to\infty$, both players effectively lose the Game of Chicken—no one wins.\footnote{However, for similar reasons commitment improves the Battle of Sexes.}
 
The dynamic described above can be characterized as a \textit{commitment race}: players have strong incentives to commit early, leading to frequent miscoordination and low expected payoffs. Of course, both players committing to Straight in the first period is not an equilibrium, but in the mixed-strategy equilibrium this outcome occurs with probability $(\tfrac{3}{4})^2 \approx 56\%$ in the Game of Chicken as $T\to\infty$. The term commitment race was originally introduced by \textcite{kokotajlo_commitment_2019} in the context of AI agents potentially extorting one another. While his treatment is highly informal, our contribution here is to formalize this idea within a game-theoretic framework.
 
\end{document}